% Derivation: Mass-Frequency Identity (m = ω in Planck units)
% Project: α-π-Helix (QNFO.RSCH.APH)
% Date: 2026-07-17
% Status: Verified — exact algebraic identity

\documentclass{article}
\usepackage{amsmath,amssymb}
\begin{document}

\section{Derivation: Mass is Angular Frequency in Planck Units}

\subsection{Fundamental Relations}

The Planck-Einstein relation connects energy and angular frequency:
\begin{equation}
E = \hbar \omega
\end{equation}

The mass-energy equivalence (special relativity):
\begin{equation}
E = m c^2
\end{equation}

\subsection{Equating}

\begin{align}
\hbar \omega &= m c^2 \\
m &= \frac{\hbar}{c^2} \omega
\end{align}

In SI units, mass and angular frequency differ by the conversion factor
$\hbar/c^2 \approx 1.173 \times 10^{-51}\ \text{kg·s}$.

\subsection{Planck Units}

In Planck units, we set:
\begin{equation}
\hbar = c = G = 1
\end{equation}

Under this normalization:
\begin{equation}
\boxed{m = \omega}
\end{equation}

Mass IS angular frequency. The "rest mass" of a particle is the
angular frequency of its internal oscillation, measured in Planck units.

\subsection{The Compton Relations}

The Compton angular frequency:
\begin{equation}
\omega_C = \frac{m c^2}{\hbar}
\end{equation}

In Planck units: $\omega_C = m$ (since $c = \hbar = 1$).

The reduced Compton wavelength:
\begin{equation}
\bar{\lambda} = \frac{\hbar}{m c}
\end{equation}

In Planck units: $\bar{\lambda} = 1/m = 1/\omega_C$.

\subsection{The Zitterbewegung Connection}

The Dirac equation gives ZB angular frequency:
\begin{equation}
\omega_{ZB} = \frac{2 m c^2}{\hbar} = 2\omega_C
\end{equation}

In Planck units: $\omega_{ZB} = 2m$.

So the ZB frequency IS (twice) the mass. The "trembling motion" is not
an extra property — it IS the mass viewed through $E = \hbar\omega$.

\subsection{α as Phase Advance}

Since $r_e = \alpha \bar{\lambda}$ and $\bar{\lambda} = c/\omega_C$:
\begin{align}
\alpha &= \frac{r_e}{\bar{\lambda}}
       = r_e \cdot \frac{\omega_C}{c}
       = \frac{r_e / c}{1 / \omega_C}
\end{align}

α is the ratio of the light-crossing time of the classical electron
radius to the Compton period. Or: the phase advance of the ZB
oscillation over the electromagnetic self-energy scale.

\subsection{Summary Table}

\begin{center}
\begin{tabular}{lcc}
\hline
Quantity & SI & Planck Units \\
\hline
Mass $m$ & $m$ (kg) & $m = \omega$ \\
Compton frequency $\omega_C$ & $mc^2/\hbar$ (rad/s) & $m$ \\
ZB frequency $\omega_{ZB}$ & $2mc^2/\hbar$ (rad/s) & $2m$ \\
Reduced Compton $\bar{\lambda}$ & $\hbar/(mc)$ (m) & $1/m$ \\
Classical radius $r_e$ & $\alpha\bar{\lambda}$ (m) & $\alpha/m$ \\
Bohr radius $a_0$ & $\bar{\lambda}/\alpha$ (m) & $1/(\alpha m)$ \\
\hline
\end{tabular}
\end{center}

\end{document}
